Diffusion image editing trades off preserving the source layout against following the target prompt. We propose Discrete Context Switching (DCS): a binary mask $M \in \{0,1\}^T$ selects which prompt (source or target) conditions each denoising timestep, producing $2^T$ candidate trajectories.

We first validate the paradigm exhaustively on Stable Diffusion~v1.5 with null-text inversion. We generate all $2^{20}\approx 10^6$ trajectories on a multi-GPU cluster for two editing tasks (cat$\to$dog, horse$\to$robot horse). Segmentation-aware metrics (background SSIM, background CLIP, foreground CLIP via a delta text embedding) score each mask, and DINOv2+UMAP+HDBSCAN split the trajectory space into 52 and 13 geometrically distinct clusters. On both tasks the winning masks activate source conditioning early and target conditioning late. The all-target trajectory $M=\mathbf{1}$ sits outside the top cluster on both tasks; running the target prompt through every step destroys the background and lowers the combined reward.

Exhaustive search costs roughly 12~GPU-hours per edit. We replace enumeration with REINFORCE policy-gradient search: a product of $N$ Bernoulli distributions, one per bit. The policy runs on the inversion-free FLUX.2-klein-base-9B flow-matching backbone ($N=14$ bits, two denoising repeats each), with entropy regularisation, variance-reduced advantages, and plateau-based early stopping. On 15 text-only background-rich editing tasks the policy converges inside 80 episodes ($\le 640$ generations), lifts the mean batch reward from $0.58$ to $0.73$ against an all-target baseline of $0.601$, and matches or beats matched-budget random search on every task. The search budget drops by two orders of magnitude relative to the $2^{14}=16{,}384$ exhaustive cost.
